\section*{Introduction}
\addcontentsline{toc}{section}{Introduction}

\subsection*{Cadre du stage}

Le stage s'est déroulé chez \acompleter{nom de l'entreprise}, au sein de \acompleter{équipe ou desk}, du
\acompleter{dates du stage}, sous la responsabilité de \acompleter{Prénom NOM, fonction}.
\acompleter{missions confiées pendant le stage, en deux ou trois phrases}.

Le présent rapport porte sur l'évaluation d'options européennes sur obligations d'État françaises (OAT). Le
contrat étudié suit les conventions de marché de la banque d'accueil : une option vanille, collatéralisée et
actualisée au taux €STR, dont le prix à terme est fixé par le financement du titre en pension livrée (repo). Les
données de marché proviennent de Bloomberg (clôture du 8 septembre 2026). Les calculs sont programmés en
Python, avec les bibliothèques NumPy, SciPy et pandas pour l'analyse numérique, statsmodels pour le filtre de
Kalman, le modèle à changement de régime et les tests statistiques, et scikit-learn pour l'analyse en
composantes principales. L'annexe~D décrit l'environnement de calcul.

\subsection*{Problématique}

Le rendement d'une OAT n'est pas le taux sans risque. Il s'en écarte d'un spread contre l'€STR qui rémunère le
risque de crédit de l'État, la liquidité du titre et des effets d'offre, et ce spread fluctue en partie
indépendamment des taux. Un pricer d'options sur obligations calibré sur les seules swaptions €STR ignore cette
source de volatilité. Or il n'existe pas de marché liquide d'options sur OAT d'où l'on pourrait la tirer : la
volatilité du spread doit être mesurée ailleurs, et il faut vérifier que les données disponibles le permettent.
Le mémoire pose donc la question suivante :

\begin{center}
\begin{minipage}{0.88\textwidth}
\itshape Faut-il modéliser le spread OAT–€STR comme un facteur de risque à part pour évaluer une option
européenne sur OAT, et peut-on en identifier les paramètres avec les données disponibles ?
\end{minipage}
\end{center}

\subsection*{Littérature et positionnement}

Le cadre retenu est le modèle gaussien à deux facteurs additifs, dit G2++, de \citet{BrigoMercurio2006}
(section~4.2), qui étend le modèle de \citet{HullWhite1990}. Le modèle recale exactement la courbe initiale, et
une option sur obligation à coupons y a un prix exact sous forme d'intégrale à une dimension (théorème~4.2.3).
À un facteur, la décomposition de \citet{Jamshidian1989} ramène cette option à une somme d'options sur
zéro-coupons. \citet{Russo} lisent le second facteur du G2++ comme un spread de crédit : le premier facteur
décrit le taux sans risque, le second l'écart entre la courbe risquée et la courbe sans risque. Ils remplacent
le prix exact par une formule fermée approchée, sous une hypothèse lognormale à «~poids gelés~», et proposent de
calibrer le facteur spread sur la seule courbe risquée.

Cette forme du taux d'actualisation se justifie dans les modèles de crédit à intensité
\citep{DuffieSingleton1999}, et \citet{MonfortRenne2014} l'appliquent aux souverains de la zone euro en séparant
crédit et liquidité ; seul le spread total entre toutefois dans le prix d'une option sur OAT.

Le travail reprend le modèle de \citet{Russo} sur des données françaises, avec trois différences. Le contrat
est celui qui se traite en banque, avec un forward fixé par le repo et non par le modèle. Le prix est calculé
exactement, à partir du théorème~4.2.3 de \citet{BrigoMercurio2006} adapté à deux courbes, et l'approximation à
poids gelés est mesurée contre lui. Enfin, l'identification des paramètres du spread est testée avant leur
estimation : la procédure de calibration du papier se révèle non identifiante, et l'estimation se fait sur les
séries historiques.

\subsection*{Hypothèses testées}

La question se décompose en cinq hypothèses, résumées au tableau~\ref{tab:hypotheses} ; la conclusion dit
pour chacune si elle est confirmée.

\begin{table}[htbp]
\centering\small
\caption{Hypothèses du mémoire et tests associés}
\label{tab:hypotheses}
\begin{tabularx}{\textwidth}{@{}l X X c@{}}
\toprule
 & Hypothèse & Test & Section \\
\midrule
H1 & La courbe OAT identifie les paramètres du facteur spread (procédure de Russo et al.) & forme de l'objectif de calibration & 5.1 \\
H2 & Les séries historiques identifient ces paramètres, sous la mesure de pricing comme sous la mesure historique & méthode des moments généralisée, régimes, filtre de Kalman, données simulées & 5, 6 \\
H3 & Le facteur spread renchérit sensiblement l'option par rapport au seul facteur taux & prix à la monnaie, décomposition de la variance, volatilités réalisées & 7.2 \\
H4 & Cet effet exige deux facteurs, et pas seulement une volatilité plus forte & comparaison à un Hull-White à un facteur équivalent & 7.3 \\
H5 & L'approximation à poids gelés de Russo et al. suffit pour le pricing & comparaison au prix exact, validé par Monte-Carlo & 7.1 \\
\bottomrule
\end{tabularx}
\end{table}

\subsection*{Enseignements de l'ENSAE mobilisés}

Trois cours de l'ENSAE fondent les choix du mémoire. Le pricing de taux s'appuie sur le cours de
\citet{Chaix2021} (actualisation au taux du collatéral, mesure forward, volatilité normale des swaptions,
Hull-White à un facteur et sa calibration sur une diagonale de swaptions) et sur celui de \citet{Hillairet2026}
(construction de la courbe, modèles de taux court à shift déterministe, probabilité forward, décomposition de
Jamshidian). L'estimation s'appuie sur le cours de \citet{MPR2025} : modèles affines gaussiens et changement de
mesure (cours~1 et~3), relation forward d'un actif stockable (cours~2), régimes de Markov (cours~4), méthode des
moments généralisée, filtre de Kalman, problème de persistance et pricing des risques de crédit et de liquidité
(cours~5). Les renvois précis, au chapitre ou à la page, figurent dans le texte.

\subsection*{Part du travail personnel}

\acompleter{rôle du maître de stage et de l'équipe : choix du sujet, données fournies, conventions de marché,
relectures ; part personnelle de l'élève dans la modélisation, la programmation et l'analyse}. La courbe de
repo, faute de cotations accessibles, est fictive et signalée comme telle à chaque usage ; la section~7 mesure
l'effet de son niveau sur les prix.

\subsection*{Plan}

La section~1 présente les données et construit les deux courbes, €STR et OAT. La section~2 expose le modèle à
deux facteurs et justifie la lecture du second facteur comme un spread. La section~3 définit le contrat, le
forward repo et le recentrage du modèle sur ce forward, puis établit le prix exact et l'approximation à poids
gelés. Les sections~4 à~6 traitent de la calibration : le facteur taux sur les swaptions, puis le facteur
spread, dont on montre d'abord que la courbe OAT ne l'identifie pas, avant de l'estimer par la méthode des
moments généralisée et par maximum de vraisemblance avec filtre de Kalman. La section~7 valide le prix exact et
présente les résultats, hypothèse par hypothèse. La conclusion reprend les réponses, les limites et les
prolongements possibles.
